% GENERATED FILE. Edit canonical lessons, then run npm run generate:books.
% canonical-source-hash: fnv1a64:8bad9ecae139d54a
% canonical-lessons: RU-C388-zakon, RU-C388-zakonnyy, RU-C388-nezakonnyy, RU-C388-narushat, RU-C388-narushenie

\chapter{The law}
\label{ch:ru-the-law}
\hlchaptermodality{388}

\begin{chapteropening}
\textbf{By the end of this chapter:} \emph{I can talk about a law, say whether something is legal or illegal, and say that someone breaks a rule.}

Complete the last lesson of chapter 388: I can talk about a law, say whether something is legal or illegal, and say that someone breaks a rule.
\end{chapteropening}

\section[\texorpdfstring{\ru{закон} (zakón)}{закон (zakón)}]{\ru{закон} (zakón) — a law}

\label{lesson:RU-C388-zakon}



\begin{quote}
Before the new one: say the Russian for to expel, then the Russian for educated.
\end{quote}

\subsection*{You'll want to know: \ru{закон}}
\textbf{\ru{закон}} (\emph{zakón}) — \textquotedblleft{}a law\textquotedblright{}.

\begin{quote}
\ru{Новый} \ru{закон} \ru{начнёт} \ru{действовать} \ru{в} \ru{январе}.
\end{quote}

\begin{quote}
\emph{Nóvyy zakón nachnyót déystvovat v yanvaré.}
\end{quote}

\begin{quote}
The new law will come into force in January.
\end{quote}

Masculine. \emph{\ru{По} \ru{закону}} — by law.

\begin{cousinweb}
\textbf{\ru{закон}} is inherited from Proto-Slavic \emph{*zakonŭ} (EtymDB), built from \emph{\ru{за}-} and \emph{\ru{кон}}, an old word for a limit or a starting line that is still heard in \emph{\ru{испокон} \ru{веков}} \textquotedblleft{}from time immemorial\textquotedblright{}.
\end{cousinweb}

\begin{grammarlens}[title={What you've built}]
One more word for telling, explaining and describing at B1.
\end{grammarlens}

\subsection*{Your turn}
\begin{itemize}
\raggedright
  \item \emph{Say it:} \emph{zakón}
  \item \emph{Say it:} \emph{zakón}, once more
  \item \emph{Say it:} a sentence of your own with \emph{zakón}
  \item \emph{Recall:} say the Russian for to expel, then the Russian for educated, then say \emph{zakón} again
\end{itemize}

\begin{tcolorbox}[breakable,colback=teal!4,colframe=teal!35!black,title={Before you move on}]
What does \ru{закон} mean? (\textquotedblleft{}A law\textquotedblright{}.) Say it once more.
\end{tcolorbox}
\section[\texorpdfstring{\ru{законный} (zakónnyy)}{законный (zakónnyy)}]{\ru{законный} (zakónnyy) — legal, lawful}

\label{lesson:RU-C388-zakonnyy}



\begin{quote}
Before the new one: say the Russian for educated, then the Russian for a law.
\end{quote}

\subsection*{You'll want to know: \ru{законный}}
\textbf{\ru{законный}} (\emph{zakónnyy}) — \textquotedblleft{}legal, lawful\textquotedblright{}.

\begin{quote}
\ru{Это} \ru{вполне} \ru{законное} \ru{требование}.
\end{quote}

\begin{quote}
\emph{Éto vpolné zakónnoye trébovaniye.}
\end{quote}

\begin{quote}
That is a perfectly lawful demand.
\end{quote}

\emph{\ru{Законно}} — legally.

\begin{cousinweb}
\textbf{\ru{законный}} is the adjective made from \emph{\ru{закон}} \textquotedblleft{}a law\textquotedblright{}, built from \emph{\ru{за}-} and old \emph{\ru{кон}} \textquotedblleft{}a limit\textquotedblright{}.
\end{cousinweb}

\begin{grammarlens}[title={What you've built}]
One more word for telling, explaining and describing at B1.
\end{grammarlens}

\subsection*{Your turn}
\begin{itemize}
\raggedright
  \item \emph{Say it:} \emph{zakónnyy}
  \item \emph{Say it:} \emph{zakónnyy}, once more
  \item \emph{Say it:} a sentence of your own with \emph{zakónnyy}
  \item \emph{Recall:} say the Russian for educated, then the Russian for a law, then say \emph{zakónnyy} again
\end{itemize}

\begin{tcolorbox}[breakable,colback=teal!4,colframe=teal!35!black,title={Before you move on}]
What does \ru{законный} mean? (\textquotedblleft{}Legal, lawful\textquotedblright{}.) Say it once more.
\end{tcolorbox}
\section[\texorpdfstring{\ru{незаконный} (nezakónnyy)}{незаконный (nezakónnyy)}]{\ru{незаконный} (nezakónnyy) — illegal, unlawful}

\label{lesson:RU-C388-nezakonnyy}



\begin{quote}
Before the new one: say the Russian for a law, then the Russian for legal.
\end{quote}

\subsection*{You'll want to know: \ru{незаконный}}
\textbf{\ru{незаконный}} (\emph{nezakónnyy}) — \textquotedblleft{}illegal, unlawful\textquotedblright{}.

\begin{quote}
\ru{Продавать} \ru{это} \ru{лекарство} \ru{без} \ru{рецепта} \ru{незаконно}.
\end{quote}

\begin{quote}
\emph{Prodavát éto lekárstvo bez retsépta nezakónno.}
\end{quote}

\begin{quote}
Selling this medicine without a prescription is illegal.
\end{quote}

\emph{\ru{Незаконно}} — illegally.

\begin{cousinweb}
\textbf{\ru{незаконный}} is \emph{\ru{не}-} \textquotedblleft{}not\textquotedblright{} in front of \emph{\ru{законный}} \textquotedblleft{}lawful\textquotedblright{}, from \emph{\ru{закон}} \textquotedblleft{}a law\textquotedblright{}.
\end{cousinweb}

\begin{grammarlens}[title={What you've built}]
One more word for telling, explaining and describing at B1.
\end{grammarlens}

\subsection*{Your turn}
\begin{itemize}
\raggedright
  \item \emph{Say it:} \emph{nezakónnyy}
  \item \emph{Say it:} \emph{nezakónnyy}, once more
  \item \emph{Say it:} a sentence of your own with \emph{nezakónnyy}
  \item \emph{Recall:} say the Russian for a law, then the Russian for legal, then say \emph{nezakónnyy} again
\end{itemize}

\begin{tcolorbox}[breakable,colback=teal!4,colframe=teal!35!black,title={Before you move on}]
What does \ru{незаконный} mean? (\textquotedblleft{}Illegal, unlawful\textquotedblright{}.) Say it once more.
\end{tcolorbox}
\section[\texorpdfstring{\ru{нарушать} (narushát)}{нарушать (narushát)}]{\ru{нарушать} (narushát) — to break (a rule), to violate}

\label{lesson:RU-C388-narushat}



\begin{quote}
Before the new one: say the Russian for legal, then the Russian for illegal.
\end{quote}

\subsection*{You'll want to know: \ru{нарушать}}
\textbf{\ru{нарушать}} (\emph{narushát}) — \textquotedblleft{}to break (a rule), to violate\textquotedblright{}.

\begin{quote}
\ru{Он} \ru{часто} \ru{нарушает} \ru{правила}.
\end{quote}

\begin{quote}
\emph{On chásto narusháyet právila.}
\end{quote}

\begin{quote}
He often breaks the rules.
\end{quote}

Imperfective; perfective \emph{\ru{нарушить}}. \emph{\ru{Нарушать} \ru{закон}, \ru{правила}, \ru{тишину}}.

\begin{cousinweb}
\textbf{\ru{нарушать}} is built from \emph{\ru{на}-} and \emph{\ru{рушить}} \textquotedblleft{}to pull down, to wreck\textquotedblright{}, an inherited Slavic verb: to break a rule is to knock part of it down.
\end{cousinweb}

\begin{grammarlens}[title={What you've built}]
One more word for telling, explaining and describing at B1.
\end{grammarlens}

\subsection*{Your turn}
\begin{itemize}
\raggedright
  \item \emph{Say it:} \emph{narushát}
  \item \emph{Say it:} \emph{narushát}, once more
  \item \emph{Say it:} a sentence of your own with \emph{narushát}
  \item \emph{Recall:} say the Russian for legal, then the Russian for illegal, then say \emph{narushát} again
\end{itemize}

\begin{tcolorbox}[breakable,colback=teal!4,colframe=teal!35!black,title={Before you move on}]
What does \ru{нарушать} mean? (\textquotedblleft{}To break (a rule), to violate\textquotedblright{}.) Say it once more.
\end{tcolorbox}
\section[\texorpdfstring{\ru{нарушение} (narushéniye)}{нарушение (narushéniye)}]{\ru{нарушение} (narushéniye) — a violation, a breach}

\label{lesson:RU-C388-narushenie}



\begin{quote}
Before the new one: say the Russian for illegal, then the Russian for to break.
\end{quote}

\subsection*{You'll want to know: \ru{нарушение}}
\textbf{\ru{нарушение}} (\emph{narushéniye}) — \textquotedblleft{}a violation, a breach\textquotedblright{}.

\begin{quote}
\ru{Это} \ru{серьёзное} \ru{нарушение} \ru{правил}.
\end{quote}

\begin{quote}
\emph{Éto seryóznoye narushéniye právil.}
\end{quote}

\begin{quote}
This is a serious breach of the rules.
\end{quote}

Neuter. \emph{\ru{Нарушение} \ru{правил}} — a breach of the rules.

\begin{cousinweb}
\textbf{\ru{нарушение}} is built on \emph{\ru{нарушить}} \textquotedblleft{}to violate\textquotedblright{}, from \emph{\ru{на}-} and \emph{\ru{рушить}} \textquotedblleft{}to pull down\textquotedblright{}, with the noun suffix \emph{-\ru{ение}}.
\end{cousinweb}

\begin{grammarlens}[title={What you've built}]
One more word for telling, explaining and describing at B1.
\end{grammarlens}

\subsection*{Your turn}
\begin{itemize}
\raggedright
  \item \emph{Say it:} \emph{narushéniye}
  \item \emph{Say it:} \emph{narushéniye}, once more
  \item \emph{Say it:} a sentence of your own with \emph{narushéniye}
  \item \emph{Recall:} say the Russian for illegal, then the Russian for to break, then say \emph{narushéniye} again
\end{itemize}

\begin{tcolorbox}[breakable,colback=teal!4,colframe=teal!35!black,title={Before you move on}]
What does \ru{нарушение} mean? (\textquotedblleft{}A violation, a breach\textquotedblright{}.) Say it once more.
\end{tcolorbox}
